\documentclass[12pt]{article}

\usepackage[a4paper, total={6in, 8in}]{geometry}
\usepackage{textcomp}

\begin{document}
\setlength{\parindent}{0pt}

\def \t {\textrightarrow}
\def \v {\vspace{1em}}
\def \bi {\begin{itemize}}
\def \ei {\end{itemize}}
\def \s[#1] {\section*{#1}}
\def \ss[#1] {\subsection*{#1}}
\def \sss[#1] {\subsubsection*{#1}}

\s[Cesare Pavese]
TUtto nasce dal neoralismo nel cinema = volonta di riproporre il reale cosi come è \t gli aristi che si ritrovano nel cinema di Luchino Visconti, in particolare "Ladri di biclette" \t attingono molto anche da Verga \\
L'esperienza della guerra, e la totale devastazione del paesaggio e dell animo delle persone \t sono stati incisivi, forti \t quindi nel cinema si riproduce il mondo cosi come è \\
Nei film si vede l umanita palpitante di tutti i giorni, e la vita quotidiana colta nella sua peculiarita \\
Moravia con "Gli indifferenti" è anche centrale \t l affondare di ogni speranza e certezza \\
Se per Gramsci l intellettuale diventa organico, per Vittorini l intellettuale deve occuparsi delle masse \t la letteratura non deve consolare, ma difendere dal dolore

\v

Qua c'è pero una vena elegiaca e intimistica \t nasce il 9 settembre 1908 (vergine) \t personaggio significativo per la nostra cultura \\
Ha soggiorno in america, dove fa incetta di testi americani \t che viene sdoganata, perche in italia non era motlo affermarata ?? \\
Gli "war poets" \t e i romanzi che caratterizzano la letteratura, come Hemigway e ??? Steneck con "La santa rossa", "Suona la campana???" \\
Ha amore per la cutlura, e la difficolta di farla passare in un area come la nostra \\
Santo Stefano Belbo \\
Lui soffre per la defezione forzata dal partito comunista?? e per l amore lontano \\
Si laurea in lettere, e conosce Tina \t si dedica a studi \\
Traduce per esempio Moby Dick \\
La sua parabola di vita si caratterizza anche per il suo insegnamento liceale \t ma la sua attivita di traduttore lo accompagna per tutta la vita \\
Lui è rappresentante del neorealismo elegiaco, come "La casa in collina", "La bella estate", e nel 1950 "La luna e il falo" che lo rende rinomato

\v

Lui si toglie la vita in un hotel \t nella sua paraobla compositiva, scrive anche "Dialoghi con Leucò??" \t qua evoca il mondo leopardiano delle operette morali, per dialogo tra figure mitologiche e dei \\
Qua si discute della morte e della vita, della sofferenza e della felicita \\
Qua complessita è quasi speculazione filosofica \\
Compone anche poesie pero \t come "Lavorare stanca", e "Verra la morte ed avra i tuoi occhi" \t sono raccolte di liriche \\
L ultima è stata oggetto di cronaca di un femminicidio

\v

Neorealsimo di tipo elegiaco è il neorealismo caratterizzato da quella vena nostalgica \t che si configura con il tema del ritorno ?? \\
I personaggi hanno marcatura introspettiva \t es. "Luna e il falo" \t devstazione di stefano belbo, paese straziato (sta nel cuore di cesare) \t macrotema di un paese devastato, e poi ricostruzione \\
Opera di Anguilla \t ragazzo che torna al suo paese, e trova un mondo diverso \t il mondo cambiato viene letto da Anguilla, come Pin che legge il mondo del "Sentiero dei nidi di ragno" \\
La ricostruzione desta in lui lo smarrimento pero \t perche lui non è piu lo stesso, ma lui non lo capisce \t le vicende non sono + le stesse, quindi non si possono replicare le cose come ?? \\
Questa ricostruzione, se è possibile il paesaggio \t non è + possibile nelle anime \t nel nostro cinema troviamo ancora le tracce del neorealismo

\v

Si constata che l intellettuale nel 900, in una societa che sta cambiando che ha ben poco da costruire sulle macerie \\
La luna e il falo non hanno + stesso significato \t hanno pari valenza \\
Falo richama le prostitute ?? \\
Nell opera c'è qundi prevalenza di monologhi interiori e riflessioni dei personaggi \t qua non c'è una provvidenza \\
Famosa trilogia
\bi
  \item "La bella estate"
  \item "La casa in collina"
  \item opera che segna la fine della parabola dell uatore
\ei

Nel 48 pubblicato "Prima che il gallo catni" \t lui muore nel 50, ed è una raccolta di due romanzi brevi: "Il carcere", e "La casa in collina"

\v

Lui e noto anche nella saggistica \t Einaudi pubblica postumo nel 68 e raccoglie gli scritti critici di Pavese \t qua ci sono:
\bi
  \item "La scoperta dell america"
  \item edgar alan poe
  \item lee master??
\ei
"On the road" \t gioventu che si approccia a torvare il suo ruolo in societa \\
In questa raccolta c e riflessione di cultura post contemporanea

\v

L opera che maggiormente accompagna la parabola dell autore è pero "Il mestiere di vivere" \t riflessioni corrosive e stroncanti \t lui è un analitico, e questi pensieri non sono datati \t non seguono sempre l ordine cronologico \\
In alcuni ci sono affermazioni lapidarie, altri sono + narrativi \t lui coglie il senso di disfacimento \\
Il mestiere di vivere è opera parcellizzata, come Zibaldone e Guicciardini \t o i pensieri di Blas Pascal, o gli aforismi di Oscar Wilde \\
Il titolo \t il mestiere di vivere è un opera che si divide tra letteratura e filosofia \t si parte da eventi qualsiasi, eventi insignicanti apparentemente \t per poi partire a riflessioni + universali \\
È il documento del mondo interiore dell artista \t sono brevi e lapidari, altri prolissi

\v

Parola chiave oltre a elegia è anche autoanalisi \t lui ha capacita di stare da solo con se stesso \t e quindi attua una delle prassi che la psiconalasi aveva annunciato, ovvero lo scavo interiore \\
Il suo bisogno di comunicazione lo porta a guardare dentro di se \t perche guardare fuori porta alla incomunicabilita \\
Solitudine, il bisogno di comunicare \t e rispetto per la donna

\v

Con cesare pavese le frontiere della nostra cultura si aprono, e la cultura italiana entra in contatto con quella americana \\
Fece scalpore la sua tragica fine \t della quale rimane la testimonianza dell autore, che si scusa \t amo una donna, che pero non corrisponde \\
Affabilita dignitosa con cui si rivolgeva alle persone \\
I cuori straziati dalla guerra \t sono i paesi e i terriotri straziati \\
Quando si parla di dolore, si parla di un risvolto nel nostro neorealismo \t che ci porta direttamente al canone del 900 \t ovvero Italo Calvino
\end{document}
